\section{Chapter 12 Conclusion and Future Scope}

\subsection{12.1 Conclusion}

This research proposed CloudSentinel AI, a context-aware cloud security framework for identifying cloud misconfigurations, constructing knowledge graphs, generating attack paths, assessing contextual risks, and providing Explainable AI-based security recommendations. The proposed framework addresses several limitations of traditional Cloud Security Posture Management (CSPM) solutions by combining graph-based reasoning with contextual analysis and AI-assisted explanations. The modular architecture, structured methodology, and evaluation plan provide a strong foundation for developing an intelligent cloud security assessment system.

\subsection{12.2 Future Scope}

Future enhancements to the proposed framework may include:

\begin{itemize}
\tightlist
\item
  Support for Microsoft Azure and Google Cloud Platform.\\
\item
  Kubernetes and container security assessment.\\
\item
  Infrastructure-as-Code (IaC) security analysis.\\
\item
  Real-time cloud monitoring and event-driven analysis.\\
\item
  Automated remediation of cloud misconfigurations.\\
\item
  Integration with SIEM and SOAR platforms.\\
\item
  Advanced machine learning models for predictive risk assessment.
\end{itemize}

\subsection{12.3 Final Remarks}

As organizations increasingly adopt cloud computing, intelligent security assessment tools become essential for protecting cloud infrastructures against evolving cyber threats. The proposed CloudSentinel AI framework aims to contribute to this domain by combining cloud security analysis, graph-based reasoning, and Explainable Artificial Intelligence into a unified and scalable solution for modern cloud environments.
